\chapter[Plan de Implantación y Puesta en Marcha]{Plan de Implantación y Puesta en Marcha}

\section*{Tipo de puesta en marcha}

OnCuidar no reemplaza un sistema anterior, porque hoy el seguimiento se lleva en
papel o de memoria. Por eso no se requiere una migración de datos ni una
operación en paralelo. Se propone una puesta en marcha gradual, que comienza con
la entrega del prototipo al equipo de Enfermería. Este equipo revisa el
contenido definitivo y decide cuándo incorporar a las primeras familias. Una
puesta en marcha gradual permite corregir el contenido y los criterios de alerta
antes de ampliar su uso, algo necesario en una aplicación de salud.

\section*{Calendario}

\begin{table}[H]
\centering
\small
\tablazebra
\caption{Calendario de implantación.}
\label{tab:implantacion}
\begin{tabularx}{\textwidth}{P{0.95}P{1.40}P{0.65}}
\toprule
\filaencabezado \encab{Etapa} & \encab{Actividades} & \encab{Fecha} \\
\midrule
Entrega del prototipo & Entrega del APK, del código fuente y de la
documentación al equipo de Enfermería. & 13 nov 2026 \\
Capacitación & Dos sesiones de taller práctico. & 16 al 18 nov 2026 \\
Revisión del contenido & Enfermería valida los textos, los materiales y los
criterios de alerta definitivos. & Por definir con Enfermería \\
Incorporación de familias & Enfermería presenta la aplicación a los primeros
cuidadores y acompaña su uso. & Por definir con Enfermería \\
\bottomrule
\end{tabularx}
\end{table}

\section*{Entorno de despliegue}

El servidor se compone de la API (Node.js y Express) y la base de datos
PostgreSQL, definidos en Docker Compose, y está destinado a ejecutarse en
contenedores Podman del servidor que la Universidad del Bío-Bío dispone para
proyectos de título. El despliegue en ese servidor está pendiente de la
activación de la cuenta solicitada; mientras tanto, el sistema funciona en un
entorno local con la misma configuración de contenedores y se probó con la
aplicación instalada en un teléfono Android. Los parámetros de operación, como
las credenciales de la base de datos, el secreto de las sesiones, la clave
maestra, el servidor de correo y la dirección pública, se entregan mediante
variables de entorno y no forman parte del repositorio. Solo se publica el
puerto de la API que asigne la universidad; la base de datos permanece en la
red interna de los contenedores.

\section*{Seguridad, respaldos, mantención y contingencia}

Los datos sensibles se cifran en el teléfono con AES-GCM antes de enviarse, de
modo que el servidor solo almacena texto cifrado y rechaza los campos clínicos
en claro. La clave de cada usuario se guarda en el almacenamiento seguro del
teléfono y una copia sellada con la clave maestra del servidor permite
recuperarla desde una sesión nueva. La clave maestra es crítica, porque sin ella
las claves de datos no se pueden abrir, por lo que debe respaldarse por separado
de la base de datos. Se recomienda programar copias periódicas de la base con
pg\_dump, almacenarlas fuera del servidor y probar su restauración antes de
incorporar familias; este respaldo aún no está automatizado. La mantención
contempla diez horas anuales por desarrolladora para corregir errores y
actualizar dependencias.

Ante las principales contingencias se definen las siguientes respuestas.

\begin{itemize}
    \item \emph{Pérdida o cambio de teléfono}: al iniciar sesión en el nuevo
    dispositivo, la aplicación solicita al servidor la clave de cifrado sellada
    y el cuidador vuelve a ver todos sus datos.
    \item \emph{Falla de conexión o del servidor}: la aplicación sigue operando
    con los datos locales, guarda las escrituras en una cola persistente y las
    reintenta con espera creciente cuando el servidor vuelve. Las escrituras
    son idempotentes, por lo que un reintento no duplica registros.
    \item \emph{Error en un contenido o en un criterio de alerta}: el contenido
    se corrige desde el backend sin publicar una nueva versión de la aplicación.
    Un criterio de alerta se corrige en el motor de reglas y se distribuye en una
    nueva versión del APK.
    \item \emph{Acceso no autorizado}: cada solicitud exige una sesión firmada
    (JWT) y el servidor filtra todo por el usuario de la sesión, de modo que un
    recurso ajeno responde como inexistente. Aun con acceso a la base de datos,
    los datos sensibles permanecen ilegibles porque están cifrados con claves
    que no se almacenan junto a ellos.
    \item \emph{Filtración de secretos o de un enlace de correo}: los enlaces de
    verificación y de recuperación son de un solo uso y vencen en una hora, y
    los secretos se rotan en las variables de entorno sin tocar el repositorio.
\end{itemize}
